The Christoffel symbols and Ricci components below are those of the derived metric of the gravitational field, $\Gamma^\rho{}_{\mu\nu} = g^\rho \cdot D_{e_\mu} g_\nu$ and $R_{\mu\nu} = -g_\mu \cdot \mathcal{R}(g_\nu)$ in the round-sphere-positive sign (Fact \claimref{synthesis.imported-gtg-dictionary}(ii) and (iv)). Write $A_{k} = a_{k}\tau^{p_{k}}$, $H_{k} := \partial_{\tau}\log A_{k} = p_{k}/\tau$, $L := \log\tau$ and $\varphi_{kj} := \partial_{j}\log A_{k} = \partial_{j}b_{k} + L\,\partial_{j}p_{k}$. For $g = d\tau^{2} - \sum_{k}A_{k}^{2}(dx^{k})^{2}$, index $0$ for $\tau$, the Christoffel symbols with a time index are $\Gamma^{0}_{kk} = A_{k}^{2}H_{k}$ and $\Gamma^{k}_{0k} = \Gamma^{k}_{k0} = H_{k}$, all others with an index $0$ vanishing (in particular $\Gamma^{0}_{0\mu} = 0$), and the spatial ones include $\Gamma^{k}_{kj} = \varphi_{kj}$. From $R_{0i} = \partial_{\mu}\Gamma^{\mu}_{0i} - \partial_{i}\Gamma^{\mu}_{0\mu} + \Gamma^{\mu}_{\mu\nu}\Gamma^{\nu}_{0i} - \Gamma^{\mu}_{i\nu}\Gamma^{\nu}_{0\mu}$: the only nonzero $\Gamma^{\mu}_{0i}$ is $\Gamma^{i}_{0i} = H_{i}$, so the first term is $\partial_{i}H_{i}$; $\Gamma^{\mu}_{0\mu} = \sum_{k}H_{k}$; the third term is $\Gamma^{\mu}_{\mu i}H_{i} = \sum_{k}\varphi_{ki}H_{i}$; and in the fourth only $\nu = \mu = k$ spatial survives, giving $\sum_{k}\Gamma^{k}_{ik}H_{k} = \sum_{k}\varphi_{ki}H_{k}$. Hence
\[
R_{0i} \;=\; \partial_{i}H_{i} - \partial_{i}\sum_{k}H_{k} + \sum_{k}\varphi_{ki}\,(H_{i} - H_{k}),
\]
and with $H_{k} = p_{k}/\tau$ and $\sum_{k}p_{k} = 1$ the second term vanishes and $\tau R_{0i} = \partial_{i}p_{i} + \sum_{k}(p_{i} - p_{k})(\partial_{i}b_{k} + L\,\partial_{i}p_{k})$. The coefficient of $L$ is $p_{i}\,\partial_{i}\sum_{k}p_{k} - \tfrac12\partial_{i}\sum_{k}p_{k}^{2} = 0$ by the Kasner relations, which gives the displayed formula; since $g_{0i} = 0$, the $0i$ component of $\mathrm{Ric} = \Lambda g$ is $R_{0i} = 0$, which is (M), for every $\Lambda$.

(i) Differentiating the Kasner relations gives $\sum_{k}\partial_{j}p_{k} = 0$ and $\sum_{k}p_{k}\partial_{j}p_{k} = 0$, so each $\partial_{j}p$ is orthogonal to $(1,1,1)$ and to $p$; these are independent because $p$ is not proportional to $(1,1,1)$ (that would need $\sum_{k}p_{k}^{2} = \tfrac13$), so $\partial_{j}p = s_{j}t$ with $t = (1,1,1) \times p = (p_{3} - p_{2}, p_{1} - p_{3}, p_{2} - p_{1}) \neq 0$ and $s$ a smooth covector. At $b = 0$, (M) reads $s_{i}t_{i} = 0$. If $s_{i} \neq 0$ on a nonempty open set $V$, then $t_{i} = 0$ on $V$, so the two exponents other than $p_{i}$ coincide there; by Proposition \claimref{sheaf.collision-census} the triples on the Kasner circle with a repeated exponent are the permutations of $(-\tfrac13, \tfrac23, \tfrac23)$ and $(0, 0, 1)$, the latter excluded by $p_{k} < 1$, a finite set, so $p$ is locally constant on $V$ and $s = 0$ there, a contradiction. Hence $s \equiv 0$, $dp = 0$, and $p$ is constant on the connected set $U$; conversely constant $p$ satisfies (M) at $b = 0$.

(ii) For $p = p(x^{1})$ the equations $i = 2, 3$ hold because no function involved depends on $x^{2}$ or $x^{3}$, and equation $1$ reads $\partial_{1}p_{1} = (p_{2} - p_{1})\partial_{1}b_{2} + (p_{3} - p_{1})\partial_{1}b_{3}$, which the displayed $b_{2}, b_{3}$ satisfy; $D > 0$ because $p_{2} = p_{3} = p_{1}$ is impossible on the circle, so the $b_{j}$ are smooth on any interval. The connection: the Lemma \claimref{dark-energy.variable-exponents-residue-pointwise-corrections-are} derives the Levi-Civita connection of the coframe $e^{0} = d\tau$, $e^{i} = A_{i}\,dx^{i}$ using only the structure coefficients $\partial_{\tau}\log A_{i}$ and $\partial_{j}\log A_{i}/A_{j}$; with $A_{i} = a_{i}\tau^{p_{i}}$ the first is still $p_{i}/\tau$, so $\omega^{i0} = (p_{i}/\tau)e^{i}$ and the residue is the pointwise Kasner boost, and the second becomes $(\partial_{j}b_{i} + L\,\partial_{j}p_{i})\,\tau^{-p_{j}}/a_{j}$, of the same order $\tau^{-p_{j}}|\log\tau|$ as in the Lemma, so its order counts for the connection and the curvature go through with the same bounds, with $|\log\tau|$ replaced by $1 + |\log\tau|$. In the wall counterterm of Proposition \claimref{dark-energy.wall-counterterm}(ii) the surviving cubic term in the boost block is a multiple of $p_{1}p_{2}p_{3}\tau^{-3}e^{1}\wedge e^{2}\wedge e^{3}$, and now $e^{1}\wedge e^{2}\wedge e^{3} = a_{1}a_{2}a_{3}\,\tau\,d^{3}x$, so its density is $8\,p_{1}p_{2}p_{3}\,a_{1}a_{2}a_{3}\,\varepsilon^{-2}d^{3}x$ at leading order.

(iii) is read off (M): equation $i$ contains $\partial_{i}b_{j}$ only, with coefficient $p_{j} - p_{i}$, and the two coefficients of one row vanish together only if $p_{1} = p_{2} = p_{3}$.

The script recomputes $R_{0i}$ independently of this derivation, as the full coordinate Ricci tensor at a point from the metric's value and its first and second derivatives over $\mathbb{Q}$, with $\log\tau$ an independent symbol value and the exponents at rational points of the circle with tangent derivatives, and finds the displayed formula at fixed rational test points (whose values it also pins), checks that the full $R_{01}$ vanishes for the one-coordinate scale functions, $\tau R_{0i} = \partial_{i}p_{i}$ at $b = 0$, and failure of the formula for derivatives off the circle (control); it confirms that $dp/du$ is parallel to $t$ in the rational parametrisation $p(u) = (-u, 1 + u, u + u^{2})/(1 + u + u^{2})$, that $t_{i} = 0$ only at collisions, and the one-coordinate solution at rational points.
